\chapter{Reproducing and diagnosing \tfunet}
\label{ch:tfunet}

\section{Initial result}
The first step was to reproduce the authors' method. I used their \tfunet{}
code without modifying the network itself, with 3 layers and 32 features in
the first layer (465{,}986 trainable parameters), and trained it on the
276$\times$600 synthetic dataset. Two settings had to differ from the paper.
Akeret et al.\ used a batch of 32 images and momentum gradient descent with a
learning rate of 0.2~\cite{akeret2017}. A batch of 32 does not fit on the
6\,GB GPU used here, so I used a batch of 4. At such a small batch the paper's
learning rate destroyed the network within a few iterations (with a batch of
one it stopped learning at iteration 4), so I used the Adam optimiser with a
learning rate of $10^{-3}$ instead.

The result was poor: a maximum \Fone{} of only 0.3879 on the test set, while my
own model (\cref{ch:hybrid}) scored about 0.98 on the same data. The obvious
conclusion would have been that the authors' architecture is simply weaker.
My first suspicion was the normalisation of the input images. Rather than
assume either explanation, I ran a set of controlled experiments.

\section{Four controlled runs}
In four runs everything was kept the same---dataset, network, optimiser and
batch size---except two settings: how the input images were normalised, and
whether class weighting was used (\cref{tab:tfunet_runs},
\cref{fig:tfunet_runs}). Class weighting gives the rare RFI pixels a larger
weight in the loss, and had been added by us to deal with the imbalance
between RFI and clean pixels; it is not part of the authors' setup.

\begin{table}[H]
  \centering
  \caption{\tfunet{} on the 276$\times$600 synthetic dataset. Only the
    normalisation and the class weighting change between runs.}
  \label{tab:tfunet_runs}
  \small
  \begin{tabular}{clcrrrl}
    \toprule
    Run & Normalisation & Class weights & Epochs & ROC AUC & max \Fone{} & Behaviour \\
    \midrule
    1 & per-image   & on  & 22 & 0.6681 & 0.3879 & frozen for 15 epochs \\
    2 & per-image   & on  & 5  & 0.5677 & 0.3064 & collapsed, stopped \\
    3 & per-image   & off & 60 & 0.8747 & 0.7191 & trained normally \\
    4 & fixed range & off & 60 & 0.9908 & 0.9317 & trained well \\
    \bottomrule
  \end{tabular}
\end{table}

\begin{figure}[H]
  \centering
  \includegraphics[width=0.75\textwidth]{tfunet_controlled_runs}
  \caption{Test \Fone{} of the four controlled \tfunet{} runs of
    \cref{tab:tfunet_runs}. The network is identical in all four.}
  \label{fig:tfunet_runs}
\end{figure}

Removing the class weights and switching to a fixed normalisation range raised
the score from 0.3879 to 0.9317 without any change to the network. Runs 3 and
4 differ only in normalisation and were trained for the same number of epochs,
so the gain from normalisation alone, $+0.213$, is a fair comparison. The
effect of class weighting is less cleanly separated: run~2 was stopped at
epoch~5 by an automatic check that detects a collapsed network, and run~1
after 22 epochs, while runs 3 and 4 trained for 60. As
\cref{sec:tfunet_v4} shows, a frozen network can eventually recover if
trained for longer, so part of the difference between runs 1--2 and run~3 is
due to training length. What the runs establish clearly is that the poor
initial score came from these two settings, not from the architecture.

\section{Cause 1: class weighting and the ReLU output}
\label{sec:relu_trap}
For every pixel, \tfunet{} produces two scores, one for ``clean'' and one for
``RFI'', and passes both through a ReLU before turning them into
probabilities. A ReLU replaces any negative value with zero. Class weighting
pushes the network strongly towards predicting RFI, which drives the
``clean'' score negative. The ReLU then sets it to exactly zero, and because
the slope of a ReLU is zero for negative inputs, no gradient flows back: the
network has no signal telling it how to change that score. It stops learning.

This is exactly what run~2 shows. It gave the same output, 0.537, for every
one of the 24{,}840{,}000 test pixels. Working backwards through the softmax,
$\mathrm{softmax}(0, c)_2 = 0.537$ gives $c = 0.148$: the ``clean'' score was
stuck at zero and the ``RFI'' score at 0.148 everywhere. The network had become
a constant.

Neither part is a problem on its own. The authors trained successfully with
the ReLU, and class weighting is a standard remedy for imbalanced data. It is
the combination that causes the network to freeze. The freeze is long rather
than permanent: on the second synthetic dataset the same network stayed frozen
for 15 epochs and then began to recover (\cref{sec:tfunet_v4}).

\section{Cause 2: per-image normalisation}
\label{sec:pernorm}
The data loader we had written scaled each image by its own minimum and
maximum. In this dataset the maximum value of an image ranges from 27 to 222,
a factor of 8.2. The same physical brightness therefore ends up at different
values in different images. \Cref{tab:pernorm} shows what this does to a noise
pixel (raw value 10) and a faint RFI pixel (raw value 20).

\begin{table}[H]
  \centering
  \caption{Normalised values of the same two pixels under per-image and
    fixed-range normalisation.}
  \label{tab:pernorm}
  \begin{tabular}{lccc}
    \toprule
    Pixel (raw value) & Per-image, quiet image & Per-image, bright image & Fixed range \\
     & (max 27) & (max 222) & \\
    \midrule
    Noise (10)     & 0.370 & 0.045 & 0.345 \\
    Faint RFI (20) & 0.741 & 0.090 & 0.520 \\
    \bottomrule
  \end{tabular}
\end{table}

With per-image scaling, noise in the quiet image (0.370) appears brighter than
real RFI in the bright image (0.090), so no single threshold separates RFI from
noise across all images. A fixed range keeps the ordering. The authors had
avoided this problem: their own launch script clips every image to the same
fixed range (from 30 to 210), but this was lost when their code was adapted to
our data. \tfunet{} has no normalisation layers inside the network, so it
cannot correct for it.

The fix was to scale all images with one fixed range, taken as the 0.5th and
99.5th percentiles of the pixel values, averaged over 40 training images
(giving $-9.69$ and $47.40$). Percentiles are used rather than the minimum and
maximum because a single extreme pixel can move the minimum or maximum a long
way. Only training images are used, so no information from the test set
enters the normalisation.

\section{The bandpass dataset}
\label{sec:tfunet_v4}
I repeated the two extreme settings on the 1024$\times$265 dataset with the
instrument bandpass. With per-image normalisation and class weights, the
network stayed frozen at a validation \Fone{} of about 0.37 for the first 15
epochs, then escaped and kept improving until the last epoch, finishing at a
test \Fone{} of 0.6815. With a fixed range and no class weights it reached
0.9483 and had converged by epoch 30. The same two settings decide the result
on the harder dataset too.

\section{\tfunet{} on real LOFAR data}
\label{sec:tfunet_lofar}
Having found the settings that make \tfunet{} work on synthetic data, I trained
it on the real LOFAR data with exactly those settings: the unmodified network,
fixed-range normalisation and no class weighting. The model was trained on the
6621 clean training images, the threshold was chosen on the 735 validation
images, and the result was measured on the 109 images labelled by the human
expert. Each setting was run with three random seeds.

Despite the extreme imbalance---one RFI pixel for every 129 clean ones---the
network did not collapse into predicting ``no RFI'' everywhere. With the same
training length as on the synthetic data (60 epochs at a learning rate of
$10^{-3}$) it reached a maximum \Fone{} of $0.4901 \pm 0.0495$. Lowering the
learning rate to $10^{-4}$ and training for 150 epochs raised this to
$0.5482 \pm 0.0139$ (\cref{tab:tfunet_lofar}).

\begin{table}[H]
  \centering
  \caption{\tfunet{} on the 109 human-labelled LOFAR test images. Mean and
    standard deviation over three seeds.}
  \label{tab:tfunet_lofar}
  \begin{tabular}{lccc}
    \toprule
    Training & max \Fone{} & pooled \Fone{} & ROC AUC \\
    \midrule
    lr $10^{-3}$, 60 epochs  & $0.4901 \pm 0.0495$ & $0.4563 \pm 0.0279$ & $0.9389 \pm 0.0045$ \\
    lr $10^{-4}$, 150 epochs & $0.5482 \pm 0.0139$ & $0.5021 \pm 0.0259$ & $0.9470 \pm 0.0074$ \\
    \bottomrule
  \end{tabular}
\end{table}

The lower learning rate also made the result much more consistent between
seeds: the standard deviation of the maximum \Fone{} fell from 0.0495 to
0.0139. However, the improvement itself is not statistically significant with
three seeds. Comparing the runs seed by seed, one of the three seeds actually
became slightly worse, and a paired $t$-test gives $t = 1.54$ with 2 degrees of
freedom, well below the 4.30 needed for significance at the 5\% level. My
first single run had suggested a gain of 0.118; the two further seeds roughly
halved it. I therefore use the lower learning rate as the better default, but
do not claim it as a proven improvement.

Using per-image normalisation instead, the maximum \Fone{} fell to 0.3039, below
even a simple threshold, which confirms on real data the conclusion of
\cref{sec:pernorm}.

\Cref{fig:tfunet_lofar} places the result among published methods. The values
for Mesarcik et al.'s U-Net and for RFI-Net are taken from their
paper~\cite{mesarcik2022}; I did not re-run those methods. Their results are
reported as maximum \Fone{}, which is why that measure is used here.
\tfunet{} clearly beats a simple $\sigma$-clipping threshold (0.4103) but stays
below \aoflagger{} (0.5698), and is 0.039 below Mesarcik et al.'s U-Net (0.5876).
Two known differences in how that U-Net was trained could account for at least
part of the gap: it was trained on small 32$\times$32 patches rather than full
images, and with a different normalisation. I have not tested which of these
matters.

\begin{figure}[H]
  \centering
  \includegraphics[width=0.8\textwidth]{tfunet_lofar_ladder}
  \caption{Maximum \Fone{} on the 109 LOFAR test images. Dark bars are my
    \tfunet{} runs (error bars: standard deviation over three seeds); grey bars
    are a $\sigma$-clipping threshold, \aoflagger{}, and the values reported by
    Mesarcik et al.~\cite{mesarcik2022}.}
  \label{fig:tfunet_lofar}
\end{figure}

\section{From synthetic to real data}
\label{sec:gap}
The same network, trained with the same settings and the same number of
gradient steps (10{,}500), scores 0.9317 on the synthetic data and 0.4901 on
LOFAR: a drop of 0.44 in \Fone{}. Even the best LOFAR result, after tuning the
learning rate, is 0.38 below the synthetic score. The reasons were set out in
\cref{sec:lofar_harder}: real RFI is far rarer, much of it is fainter than the
noise, and the training labels themselves are often wrong. A method that
performs well on a synthetic benchmark can therefore perform much worse on
real observations, and synthetic results alone should be read with care.
